% ============================================================
% Nature Portfolio LaTeX Template
% Compatible with: Nature, Nature Materials, Nature Chemistry,
%   Nature Energy, Nature Catalysis, Nature Communications, etc.
% ============================================================
\documentclass[9pt,twoside,lineno]{article}

% ---- Packages ----
\usepackage[utf8]{inputenc}
\usepackage[T1]{fontenc}
\usepackage{amsmath,amssymb}
\usepackage{graphicx}
\usepackage{xcolor}
\usepackage{hyperref}
\usepackage{natbib}
\usepackage{geometry}
\usepackage{setspace}
\usepackage{titlesec}
\usepackage{caption}
\usepackage{booktabs}
\usepackage{siunitx}
\usepackage{chemformula}  % for chemical formulas: \ch{H2O}
\usepackage{lineno}

\geometry{
  a4paper,
  top=2.5cm, bottom=2.5cm,
  left=2.5cm, right=2.5cm
}

\hypersetup{
  colorlinks=true,
  linkcolor=blue,
  citecolor=blue,
  urlcolor=blue
}

% Line numbers (required for submission)
\linenumbers

% ---- Journal/Article Info (fill in) ----
\newcommand{\journalname}{Nature [Sub-journal]}
\newcommand{\articletype}{Article}  % Article / Letter / Brief Communication

% ============================================================
\begin{document}

% ---- Title ----
\begin{center}
  {\LARGE\bfseries [Title: concise, informative, max ~15 words]}

  \vspace{1em}

  % Authors (use superscripts for affiliations)
  Author One$^{1}$, Author Two$^{2}$, Corresponding Author$^{1,*}$

  \vspace{0.5em}

  \footnotesize
  $^{1}$Department, Institution, City, Country \\
  $^{2}$Department, Institution, City, Country \\
  $^{*}$Correspondence: email@institution.edu

  \vspace{0.5em}
  \normalsize
  \textbf{Keywords:} keyword1; keyword2; keyword3; keyword4; keyword5
\end{center}

\hrule
\vspace{1em}

% ============================================================
% ABSTRACT (150–200 words depending on journal)
% Nature style: single paragraph, no citations, no undefined abbreviations
% ============================================================
\noindent\textbf{Abstract}

\noindent
[One-paragraph abstract. Begin with the broad context (1–2 sentences). 
State the problem or gap. Summarize the approach and key findings. 
End with the significance/impact. No references, no abbreviations undefined 
in the abstract itself. Target: 150 words for most Nature journals.]

\vspace{1em}
\hrule
\vspace{1em}

% ============================================================
% INTRODUCTION (typically 4–6 paragraphs, ~800 words)
% Nature style: end with a clear statement of what this paper reports
% ============================================================
\section*{Introduction}

[Paragraph 1: Broad context and importance of the field.]

[Paragraph 2: More specific background — what is known.]

[Paragraph 3: The gap, challenge, or open question.]

[Paragraph 4: Brief summary of your approach and key results.]

[Final sentence: "Here we report / show / demonstrate..."]

% ============================================================
% RESULTS (core section — use descriptive subsection titles)
% Each subsection should tell a story: observation → interpretation
% ============================================================
\section*{Results}

\subsection*{[Subsection 1: e.g., Synthesis and structural characterization]}

[Describe experiment/computation. Report observations with data. 
Reference figures inline: (Fig.~\ref{fig:1}a). 
Nature style: present data first, then interpret — avoid opinion without data.]

\subsection*{[Subsection 2: e.g., Mechanistic investigation]}

[...]

\subsection*{[Subsection 3: e.g., Performance evaluation]}

[...]

% ============================================================
% DISCUSSION (~400–600 words)
% Interpret results, compare with literature, state limitations, outlook
% ============================================================
\section*{Discussion}

[Paragraph 1: Summarize key findings and their significance.]

[Paragraph 2: Comparison with existing literature — how your work advances the field.]

[Paragraph 3: Mechanistic or theoretical interpretation.]

[Paragraph 4: Limitations and future directions.]

[Concluding sentence: State the broader impact.]

% ============================================================
% METHODS
% Location depends on journal:
%   - Nature, Nature Materials, etc.: AFTER references
%   - Nature Communications: within main text
% Subsections should mirror Results subsections
% ============================================================
\section*{Methods}

\subsection*{Materials}
[Chemicals/materials used, purity, source.]

\subsection*{Synthesis / Sample Preparation}
[Step-by-step synthesis protocol. Include quantities, temperatures, times.]

\subsection*{Characterization}
[Instruments, settings, standards used.]

\subsection*{[Computational Methods / DFT / MD]}
[Software, functional, basis set, pseudopotentials, k-points, cutoff, etc.]

\subsection*{Data Analysis}
[Statistical methods, software, error analysis.]

% ============================================================
% REFERENCES
% Nature style: numbered, order of appearance, Vancouver-style
% Format: Author(s). Title. Journal Vol, pages (Year).
% ============================================================
\bibliographystyle{naturemag}  % use naturemag.bst
\bibliography{references}      % references.bib file

% ============================================================
% FIGURE LEGENDS (after references for most Nature journals)
% ============================================================
\section*{Figure Legends}

\noindent\textbf{Figure 1 | [Short title describing the figure.]}
[Detailed legend. Each panel labeled a, b, c... described in order.
Include scale bars, error bars (mean ± s.d., n = X independent experiments),
statistical test used. All abbreviations defined.]

\noindent\textbf{Figure 2 | [Short title.]}
[...]

% ============================================================
% EXTENDED DATA (Nature/Nature sub-journals; up to 10 items)
% ============================================================
% \section*{Extended Data}
% Extended Data Fig. 1 | ...

% ============================================================
% SUPPLEMENTARY INFORMATION (reference in main text as ref. X)
% ============================================================
% Supplementary Information is submitted as a separate file.

% ============================================================
% END MATTER
% ============================================================
\section*{Data Availability}
[State where data is deposited. E.g., "Source data for Figs. 1–4 are provided 
with this paper. All other data are available from the corresponding author 
upon reasonable request."]

\section*{Code Availability}
[If applicable: "Code used for analysis is available at [GitHub URL]."]

\section*{Acknowledgements}
[Funding sources (grant numbers), technical assistance, facilities used.]

\section*{Author Contributions}
[Use CRediT taxonomy. E.g., "X.X. conceived the study. X.X. and Y.Y. performed experiments. All authors discussed results and wrote the manuscript."]

\section*{Competing Interests}
[Declare or state: "The authors declare no competing interests."]

\section*{Additional Information}
Correspondence and requests for materials should be addressed to [Name] ([email]).

% ============================================================
% FIGURES (can be at end for submission, or separate files)
% ============================================================
\newpage
\begin{figure}[h!]
  \centering
  \includegraphics[width=\textwidth]{fig1}
  \caption{\textbf{Figure 1 | Title.} (a) ... (b) ... 
  Scale bar, X nm. Error bars represent mean ± s.d. ($n = X$).}
  \label{fig:1}
\end{figure}

\end{document}
